\documentclass[UTF8,a4paper,12pt]{ctexart}
\ctexset{
  section/format = {\heiti\zihao{3}},
  subsection/format = {\kaishu\zihao{3}\bfseries},
}
\usepackage{geometry}
\geometry{top=37mm,bottom=35mm,left=28mm,right=26mm}
\usepackage{indentfirst}
\setlength{\parindent}{2em}
\usepackage{longtable,booktabs}
\usepackage{fancyhdr}
\pagestyle{fancy}
\fancyhf{}
\fancyhead[C]{\kaishu\zihao{6} 中普咨询 | 内部报告 | 机密}
\fancyfoot[C]{\thepage}
\usepackage{hyperref}
\hypersetup{colorlinks=true,linkcolor=black}
\setCJKmainfont{仿宋}
\setCJKsansfont{黑体}
\setCJKmonofont{楷体}
\title{个人闲置算力出租市场调研报告}
\date{2026年6月21日}
\begin{document}
\maketitle
\tableofcontents\newpage
\subsection{目录}\label{ux76eeux5f55}

\begin{center}\rule{0.5\linewidth}{0.5pt}\end{center}

\pagebreak

\subsection{1. 执行摘要}\label{ux6267ux884cux6458ux8981}

\subsubsection{核心发现}\label{ux6838ux5fc3ux53d1ux73b0}

本报告基于对全球个人闲置算力出租市场的系统调研，覆盖平台生态、经济模型、政策法规、竞争格局和战略机遇五个维度，形成以下核心判断：

\textbf{市场判断：}
全球去中心化算力共享市场正处于爆发前夜。2024年全球GPU云市场规模约\$320亿，而去中心化/P2P算力共享细分市场约\$15-25亿，渗透率不足8\%。在AI大模型训练推理需求激增（年增速\textgreater40\%）和GPU供给短缺的双重推动下，预计2030年该细分市场将达\$80-150亿（CAGR
30-45\%）。中国市场受''东数西算''国家战略和算力网络建设驱动，个人算力共享正在从灰色地带走向政策鼓励的规范化发展轨道。

\textbf{技术判断：}
算力共享的技术栈正在快速成熟------容器化部署（Docker/K8s）、GPU虚拟化（MIG/vGPU）、区块链结算、去中心化调度，四大技术支柱已基本就绪。关键瓶颈在于网络延迟（对分布式训练影响显著）和可信执行环境（TEE）的安全认证。边缘计算与算力共享的融合是下一阶段的技术演进方向。

\textbf{竞争判断：}
全球形成''四类玩家''格局------（1）中心化GPU云平台（Vast.ai、RunPod），以P2P模式连接供需；（2）区块链去中心化算力网络（Akash、Render
Network、io.net），以代币经济激励供给；（3）消费级共享应用（Salad、Golem），聚焦游戏PC闲置资源；（4）中国本土平台（趋动云、矩池云、阿里云GPU共享），受政策监管影响更为谨慎。当前平台方抽成约15-25\%，供给端竞争激烈，单卡收益呈下降趋势。

\textbf{战略建议：}
对个人算力出租者，建议采取''平台组合+时段优化''策略------优先选择Vast.ai（全球流动性最好）或RunPod（更稳定），长期部署选高价时段（北美工作日白天）。对机构和投资者，建议关注三个战略机会点：（1）算力聚合运营商（帮助个人管理多平台）；（2）边缘算力节点基础设施（结合5G基站下沉）；（3）合规化算力交易平台（中国政策窗口）。

\subsubsection{关键数字}\label{ux5173ux952eux6570ux5b57}

{\def\LTcaptype{none} % do not increment counter
\begin{longtable}[]{@{}ll@{}}
\toprule\noalign{}
指标 & 数值 \\
\midrule\noalign{}
\endhead
\bottomrule\noalign{}
\endlastfoot
2024年全球GPU云市场规模 & ≈\$320亿 \\
去中心化算力共享市场规模 & ≈\$15-25亿 \\
2030年去中心化算力共享市场（预测） & ≈\$80-150亿 \\
Vast.ai在线GPU数量 & \textgreater20,000张 \\
RTX 4090月均出租收益（全球均价） & \$150-350/月 \\
RTX 3090月均出租收益 & \$100-200/月 \\
A100 80GB月均出租收益 & \$500-1,200/月 \\
平台方抽成比例 & 15-25\% \\
中国个人GPU持有量（游戏级及以上） & ≈800-1,200万张 \\
个人算力出租回本周期 & 12-24个月 \\
全球算力共享活跃供给节点 & ≈15-30万个 \\
\end{longtable}
}

\begin{center}\rule{0.5\linewidth}{0.5pt}\end{center}

\pagebreak

\subsection{2.
行业定义与核心特性}\label{ux884cux4e1aux5b9aux4e49ux4e0eux6838ux5fc3ux7279ux6027}

\subsubsection{2.1 定义}\label{ux5b9aux4e49}

个人闲置算力出租（Personal Idle Compute Power
Rental）是指个人或小型组织将其拥有的GPU/CPU计算资源在空闲时段通过第三方平台或点对点（P2P）协议出租给其他用户，以获取经济回报的经济行为。这是共享经济（Sharing
Economy）在算力领域的延伸，也被称为算力共享（Compute
Sharing）、分布式GPU云（Distributed GPU
Cloud）或去中心化算力市场（Decentralized Compute Marketplace）。

\textbf{核心特征：}

\begin{enumerate}
\def\labelenumi{\arabic{enumi}.}
\tightlist
\item
  \textbf{供给端去中心化：}
  供给方为分散的个人或小规模数据中心，而非AWS/Azure/GCP等集中式云厂商
\item
  \textbf{价格市场化：}
  价格由供需关系实时决定，波动幅度显著大于传统云服务
\item
  \textbf{轻运营模式：} 平台方不持有硬件资产，仅提供匹配、调度和结算服务
\item
  \textbf{全球分布：} 算力节点散布全球，时区差异形成自然的供需互补
\end{enumerate}

\subsubsection{2.2
分类与细分领域}\label{ux5206ux7c7bux4e0eux7ec6ux5206ux9886ux57df}

\paragraph{2.2.1
按平台架构分类}\label{ux6309ux5e73ux53f0ux67b6ux6784ux5206ux7c7b}

{\def\LTcaptype{none} % do not increment counter
\begin{longtable}[]{@{}
  >{\raggedright\arraybackslash}p{(\linewidth - 6\tabcolsep) * \real{0.1818}}
  >{\raggedright\arraybackslash}p{(\linewidth - 6\tabcolsep) * \real{0.2727}}
  >{\raggedright\arraybackslash}p{(\linewidth - 6\tabcolsep) * \real{0.2727}}
  >{\raggedright\arraybackslash}p{(\linewidth - 6\tabcolsep) * \real{0.2727}}@{}}
\toprule\noalign{}
\begin{minipage}[b]{\linewidth}\raggedright
类型
\end{minipage} & \begin{minipage}[b]{\linewidth}\raggedright
代表平台
\end{minipage} & \begin{minipage}[b]{\linewidth}\raggedright
核心机制
\end{minipage} & \begin{minipage}[b]{\linewidth}\raggedright
适用场景
\end{minipage} \\
\midrule\noalign{}
\endhead
\bottomrule\noalign{}
\endlastfoot
\textbf{中心化P2P市场} & Vast.ai, RunPod & 平台集中匹配+调度，法币结算 &
AI训练/推理、渲染 \\
\textbf{区块链去中心化网络} & Akash, Render, io.net &
智能合约调度，代币结算 & Web3应用、渲染 \\
\textbf{消费级共享应用} & Salad, Golem, NiceHash &
客户端后台运行，自动匹配任务 & 挖矿、科学计算 \\
\textbf{企业级算力交易} & CoreWeave (早期), FluidStack &
B2B协议，长期合约 & 大规模AI训练 \\
\end{longtable}
}

\paragraph{2.2.2
按硬件供给层级分类}\label{ux6309ux786cux4ef6ux4f9bux7ed9ux5c42ux7ea7ux5206ux7c7b}

{\def\LTcaptype{none} % do not increment counter
\begin{longtable}[]{@{}
  >{\raggedright\arraybackslash}p{(\linewidth - 6\tabcolsep) * \real{0.1622}}
  >{\raggedright\arraybackslash}p{(\linewidth - 6\tabcolsep) * \real{0.2432}}
  >{\raggedright\arraybackslash}p{(\linewidth - 6\tabcolsep) * \real{0.2973}}
  >{\raggedright\arraybackslash}p{(\linewidth - 6\tabcolsep) * \real{0.2973}}@{}}
\toprule\noalign{}
\begin{minipage}[b]{\linewidth}\raggedright
层级
\end{minipage} & \begin{minipage}[b]{\linewidth}\raggedright
典型硬件
\end{minipage} & \begin{minipage}[b]{\linewidth}\raggedright
月租金范围
\end{minipage} & \begin{minipage}[b]{\linewidth}\raggedright
主要租用方
\end{minipage} \\
\midrule\noalign{}
\endhead
\bottomrule\noalign{}
\endlastfoot
\textbf{消费级} & RTX 3060/4060, GTX 1660 & \$30-80/月 &
小型AI微调、游戏渲染 \\
\textbf{高端消费级} & RTX 4090, RTX 3090 & \$150-400/月 &
LLM推理、Stable Diffusion \\
\textbf{专业级} & RTX A6000, A5000 & \$250-600/月 &
专业渲染、中小训练 \\
\textbf{数据中心级} & A100, H100, B200 & \$500-3,000/月 &
大规模训练、HPC \\
\end{longtable}
}

\subsubsection{2.3
核心参数对比表}\label{ux6838ux5fc3ux53c2ux6570ux5bf9ux6bd4ux8868}

{\def\LTcaptype{none} % do not increment counter
\begin{longtable}[]{@{}
  >{\raggedright\arraybackslash}p{(\linewidth - 8\tabcolsep) * \real{0.0682}}
  >{\raggedright\arraybackslash}p{(\linewidth - 8\tabcolsep) * \real{0.2386}}
  >{\raggedright\arraybackslash}p{(\linewidth - 8\tabcolsep) * \real{0.2386}}
  >{\raggedright\arraybackslash}p{(\linewidth - 8\tabcolsep) * \real{0.2386}}
  >{\raggedright\arraybackslash}p{(\linewidth - 8\tabcolsep) * \real{0.2159}}@{}}
\toprule\noalign{}
\begin{minipage}[b]{\linewidth}\raggedright
维度
\end{minipage} & \begin{minipage}[b]{\linewidth}\raggedright
传统云计算 (AWS/GCP)
\end{minipage} & \begin{minipage}[b]{\linewidth}\raggedright
中心化P2P (Vast.ai)
\end{minipage} & \begin{minipage}[b]{\linewidth}\raggedright
去中心化网络 (Akash)
\end{minipage} & \begin{minipage}[b]{\linewidth}\raggedright
消费级共享 (Salad)
\end{minipage} \\
\midrule\noalign{}
\endhead
\bottomrule\noalign{}
\endlastfoot
\textbf{价格 (A100/hr)} & \$3.00-5.00 & \$0.80-1.50 & \$0.60-1.20 &
N/A \\
\textbf{价格 (RTX 4090/hr)} & 不提供 & \$0.30-0.60 & \$0.25-0.50 &
以积分结算 \\
\textbf{最低起租} & 按需/预留实例 & \$5账户充值 & ≈\$1 (AKT代币) &
免费客户端 \\
\textbf{可靠性 SLA} & 99.95\%+ & 无SLA (约95-98\%) & 无SLA (约90-95\%) &
无SLA \\
\textbf{部署速度} & 分钟级 & 秒级 & 分钟级 & 自动分配 \\
\textbf{数据安全} & 高 (VPC隔离) & 中低 (依赖主机) & 中 (TEE可选) &
低 \\
\textbf{网络延迟} & 低 (专线可选) & 不确定 (1-200ms) & 不确定 &
不确定 \\
\textbf{全球节点} & 30+区域 & 40+数据中心 & 100+提供商 & 190+国家 \\
\textbf{结算货币} & 法币/信用 & 法币 (信用卡) & AKT等代币 &
积分/礼品卡 \\
\textbf{供给方资质} & 企业认证 & 个人+KYC简化 & 无许可 &
消费级自动验证 \\
\end{longtable}
}

\begin{center}\rule{0.5\linewidth}{0.5pt}\end{center}

\pagebreak

\subsection{3.
产业链与价值链分析}\label{ux4ea7ux4e1aux94feux4e0eux4ef7ux503cux94feux5206ux6790}

\subsubsection{3.1 产业链全景}\label{ux4ea7ux4e1aux94feux5168ux666f}

\begin{verbatim}
┌─────────────────────────────────────────────────────────────────┐
│                        算力共享产业链                              │
├──────────┬──────────┬──────────┬──────────┬──────────┬──────────┤
│ 硬件供给  │ 平台/协议 │ 算力聚合  │ 调度分发  │ 终端需求  │ 配套服务  │
├──────────┼──────────┼──────────┼──────────┼──────────┼──────────┤
│ GPU厂商   │ Vast.ai  │ 算力运营商 │ 容器引擎  │ AI创业公司 │ 保险/担保  │
│ (NVIDIA) │ RunPod   │ (汇集个人 │ (Docker)  │ 科研机构  │ 硬件融资  │
│ AMD      │ Akash    │  节点出售) │ K8s      │ 渲染工作室 │ 电力优化  │
│ Intel    │ Render   │ FluidStack │ Slurm    │ Web3项目  │ 散热方案  │
│          │ io.net   │           │          │ 量化交易  │ 合规咨询  │
│ 个人持有者│ Salad    │           │ 区块链    │ 个人开发者 │ 运维托管  │
│ 小型DC   │ Golem    │           │ 智能合约   │          │          │
│ 网吧     │ NiceHash │           │          │          │          │
│ 矿场转型  │          │           │          │          │          │
└──────────┴──────────┴──────────┴──────────┴──────────┴──────────┘
\end{verbatim}

\subsubsection{3.2
价值分配模型}\label{ux4ef7ux503cux5206ux914dux6a21ux578b}

以一张RTX 4090在Vast.ai平台上的典型出租为例（假设月租金\$300）：

{\def\LTcaptype{none} % do not increment counter
\begin{longtable}[]{@{}
  >{\raggedright\arraybackslash}p{(\linewidth - 6\tabcolsep) * \real{0.2222}}
  >{\raggedright\arraybackslash}p{(\linewidth - 6\tabcolsep) * \real{0.3333}}
  >{\raggedright\arraybackslash}p{(\linewidth - 6\tabcolsep) * \real{0.2222}}
  >{\raggedright\arraybackslash}p{(\linewidth - 6\tabcolsep) * \real{0.2222}}@{}}
\toprule\noalign{}
\begin{minipage}[b]{\linewidth}\raggedright
环节
\end{minipage} & \begin{minipage}[b]{\linewidth}\raggedright
分配金额
\end{minipage} & \begin{minipage}[b]{\linewidth}\raggedright
占比
\end{minipage} & \begin{minipage}[b]{\linewidth}\raggedright
说明
\end{minipage} \\
\midrule\noalign{}
\endhead
\bottomrule\noalign{}
\endlastfoot
GPU持有者（个人） & \$225-255 & 75-85\% & 扣除平台费后的净收入 \\
平台方 & \$45-75 & 15-25\% & Vast.ai抽佣约15\%，部分平台更高 \\
电力成本（月耗电约150-200kWh） & \$15-30 & 5-10\% &
按\$0.10-0.15/kWh计算 \\
硬件折旧（3年直线折旧） & \$45-60 & 15-20\% & RTX 4090 ≈\$1,600/36月 \\
净收益（税前） & \$135-180 & 45-60\% & 扣除所有成本后的纯收益 \\
\end{longtable}
}

\textbf{关键发现：} 在\$300/月的乐观场景下，RTX
4090约10-14个月可收回硬件投资。但实际平均出租率约50-70\%，月均收益在\$150-250之间更为常见，回本周期延长到16-24个月。

\subsubsection{3.3
价值驱动因素}\label{ux4ef7ux503cux9a71ux52a8ux56e0ux7d20}

\begin{itemize}
\tightlist
\item
  \textbf{AI算力供需缺口：} NVIDIA
  GPU供不应求，H100/B200交付周期长达数月，溢出需求流向P2P市场
\item
  \textbf{价格优势：} P2P算力比AWS/GCP同类GPU便宜50-80\%
\item
  \textbf{时区套利：} 亚洲节点在北美白天时段出租，实现24小时利用
\item
  \textbf{尾部需求满足：}
  小型团队和个人开发者无法获得云计算的最低承诺用量
\end{itemize}

\begin{center}\rule{0.5\linewidth}{0.5pt}\end{center}

\pagebreak

\subsection{4.
市场规模与增长预测}\label{ux5e02ux573aux89c4ux6a21ux4e0eux589eux957fux9884ux6d4b}

\subsubsection{4.1 全球市场}\label{ux5168ux7403ux5e02ux573a}

{\def\LTcaptype{none} % do not increment counter
\begin{longtable}[]{@{}lllll@{}}
\toprule\noalign{}
年份 & GPU云市场总规模 & 去中心化/P2P细分 & 渗透率 & 活跃供给节点数 \\
\midrule\noalign{}
\endhead
\bottomrule\noalign{}
\endlastfoot
2022 & ≈\$150亿 & ≈\$5-8亿 & 3-5\% & ≈5-10万 \\
2023 & ≈\$220亿 & ≈\$10-15亿 & 5-7\% & ≈8-15万 \\
2024 & ≈\$320亿 & ≈\$15-25亿 & 5-8\% & ≈15-30万 \\
2025E & ≈\$460亿 & ≈\$30-50亿 & 7-11\% & ≈25-50万 \\
2027E & ≈\$800亿 & ≈\$60-100亿 & 8-13\% & ≈50-100万 \\
2030E & ≈\$1,500亿 & ≈\$80-150亿 & 5-10\% & ≈80-200万 \\
\end{longtable}
}

\emph{数据来源：Grand View Research, Fortune Business Insights, Messari,
中普咨询整理估算}

\subsubsection{4.2
市场增长驱动力}\label{ux5e02ux573aux589eux957fux9a71ux52a8ux529b}

\begin{enumerate}
\def\labelenumi{\arabic{enumi}.}
\tightlist
\item
  \textbf{AI大模型军备竞赛：}
  GPT-5级模型训练需万卡集群，头部企业外溢需求巨大
\item
  \textbf{推理需求爆发：} Agent应用普及（如Hermes
  Agent类产品），推理算力需求年增200\%+
\item
  \textbf{GPU产能瓶颈：}
  NVIDIA年产能无法满足需求，H200/B300交付周期达6-12个月
\item
  \textbf{边缘计算崛起：} 5G+AI推动边缘推理，个人节点地理位置优势显现
\item
  \textbf{Web3+AI融合：}
  去中心化AI训练和推理成为新范式，代币激励加速节点增长
\end{enumerate}

\subsubsection{4.3
中国市场特殊性分析}\label{ux4e2dux56fdux5e02ux573aux7279ux6b8aux6027ux5206ux6790}

中国个人算力出租市场面临独特的机遇与约束：

\textbf{机遇：} -
``东数西算''工程总投资超¥4,000亿，算力网络建设纳入国家战略 -
中国游戏GPU保有量全球第一（Steam硬件调查中简体中文用户占比\textgreater30\%）
- 大量加密货币矿场转型需求，闲置A卡/N卡存量可观 -
AI创业公司数量激增，中小团队算力缺口明显

\textbf{约束：} - 跨境数据流动受《数据安全法》《个人信息保护法》严格限制
- 加密货币相关政策不确定性影响区块链算力平台 -
家用宽带商业使用受运营商服务条款限制 - 备案与许可证要求提高合规门槛

{\def\LTcaptype{none} % do not increment counter
\begin{longtable}[]{@{}ll@{}}
\toprule\noalign{}
指标 & 中国市场估算 \\
\midrule\noalign{}
\endhead
\bottomrule\noalign{}
\endlastfoot
中国GPU云市场规模（2024年） & ≈¥450-600亿（\$60-80亿） \\
个人闲置GPU可激活总量 & ≈300-500万张（游戏卡+矿卡） \\
去中心化/P2P算力占比 & \textless2\%（远低于全球5-8\%） \\
合规化率 & \textless30\% \\
主流通证结算平台使用率 & \textless5\%（受政策限制） \\
\end{longtable}
}

\begin{center}\rule{0.5\linewidth}{0.5pt}\end{center}

\pagebreak

\subsection{5.
需求端深度分析}\label{ux9700ux6c42ux7aefux6df1ux5ea6ux5206ux6790}

\subsubsection{5.1
核心需求场景}\label{ux6838ux5fc3ux9700ux6c42ux573aux666f}

\paragraph{5.1.1
AI模型训练与微调（占比约45\%）}\label{aiux6a21ux578bux8badux7ec3ux4e0eux5faeux8c03ux5360ux6bd4ux7ea645}

{\def\LTcaptype{none} % do not increment counter
\begin{longtable}[]{@{}llll@{}}
\toprule\noalign{}
租用方画像 & 典型需求 & 租用时长 & 价格敏感度 \\
\midrule\noalign{}
\endhead
\bottomrule\noalign{}
\endlastfoot
AI创业公司（种子/A轮） & LoRA微调、小型训练 & 数天-数周 & 极高 \\
独立研究者/博士生 & 实验验证、论文复现 & 数小时-数天 & 极高 \\
中小企业AI团队 & 模型蒸馏、评估 & 数周-数月 & 高 \\
Kaggle/竞赛参与者 & 短期高强度训练 & 数天 & 中 \\
\end{longtable}
}

\textbf{典型案例：}
某AI创业者使用Vast.ai租用8×A100集群进行7B模型微调，总花费\$280（vs
AWS等效\$1,200+），节省77\%。

\paragraph{5.1.2
AI推理与Agent运行（占比约25\%）}\label{aiux63a8ux7406ux4e0eagentux8fd0ux884cux5360ux6bd4ux7ea625}

\begin{itemize}
\tightlist
\item
  \textbf{AI Agent应用}（如Hermes
  Agent、AutoGPT等）需要7×24小时GPU推理资源
\item
  \textbf{图像/视频生成服务}（Stable
  Diffusion、可灵AI等）请求量波动大，弹性需求显著
\item
  \textbf{语音识别/合成}（Whisper、TTS）对延迟敏感但计算量中等
\item
  \textbf{趋势：} 推理需求正在超越训练需求，成为算力增长的最大驱动力
\end{itemize}

\paragraph{5.1.3
渲染与内容制作（占比约20\%）}\label{ux6e32ux67d3ux4e0eux5185ux5bb9ux5236ux4f5cux5360ux6bd4ux7ea620}

{\def\LTcaptype{none} % do not increment counter
\begin{longtable}[]{@{}
  >{\raggedright\arraybackslash}p{(\linewidth - 6\tabcolsep) * \real{0.2571}}
  >{\raggedright\arraybackslash}p{(\linewidth - 6\tabcolsep) * \real{0.2571}}
  >{\raggedright\arraybackslash}p{(\linewidth - 6\tabcolsep) * \real{0.2286}}
  >{\raggedright\arraybackslash}p{(\linewidth - 6\tabcolsep) * \real{0.2571}}@{}}
\toprule\noalign{}
\begin{minipage}[b]{\linewidth}\raggedright
渲染类型
\end{minipage} & \begin{minipage}[b]{\linewidth}\raggedright
典型软件
\end{minipage} & \begin{minipage}[b]{\linewidth}\raggedright
GPU偏好
\end{minipage} & \begin{minipage}[b]{\linewidth}\raggedright
需求特征
\end{minipage} \\
\midrule\noalign{}
\endhead
\bottomrule\noalign{}
\endlastfoot
3D动画渲染 & Blender, Octane, Redshift & RTX 4090/3090多卡 &
批处理，时效要求中等 \\
建筑可视化 & V-Ray, Enscape & RTX A系列 & 周期性（项目制） \\
视频后期 & DaVinci Resolve & RTX 4090 & 短期高爆发 \\
AI视频生成 & 可灵、Sora类 & H100/A100 & 24×7持续 \\
\end{longtable}
}

\paragraph{5.1.4
科学计算与Web3（占比约10\%）}\label{ux79d1ux5b66ux8ba1ux7b97ux4e0eweb3ux5360ux6bd4ux7ea610}

\begin{itemize}
\tightlist
\item
  分子动力学模拟（Folding@home类分布式计算）
\item
  加密货币挖矿（ETH转向PoS后GPU挖矿退潮，但AI训练替代需求上升）
\item
  零知识证明（ZK-Proof）生成
\item
  区块链节点验证
\end{itemize}

\subsubsection{5.2
需求方决策因素}\label{ux9700ux6c42ux65b9ux51b3ux7b56ux56e0ux7d20}

\begin{verbatim}
选择P2P算力 vs 传统云的主要考量（按重要性排序）：

1. ════════════════════ 价格（权重：40%）
   P2P便宜50-80%是最核心吸引力

2. ════════════════ 可用性（权重：25%）
   传统云A100/H100经常缺货，P2P可即时获取

3. ════════════ 灵活性（权重：15%）
   按时计费、无最低承诺、可随时切换节点

4. ════════ 免审核（权重：10%）
   个人开发者无需企业认证或信用审核

5. ═══ 特定配置（权重：10%）
   需要特定GPU型号或驱动版本时P2P选择更多
\end{verbatim}

\begin{center}\rule{0.5\linewidth}{0.5pt}\end{center}

\pagebreak

\subsection{6.
竞争格局与情报}\label{ux7adeux4e89ux683cux5c40ux4e0eux60c5ux62a5}

\subsubsection{6.1
全球主要平台对比}\label{ux5168ux7403ux4e3bux8981ux5e73ux53f0ux5bf9ux6bd4}

\paragraph{6.1.1 Vast.ai --- 全球最大GPU
P2P市场}\label{vast.ai-ux5168ux7403ux6700ux5927gpu-p2pux5e02ux573a}

{\def\LTcaptype{none} % do not increment counter
\begin{longtable}[]{@{}ll@{}}
\toprule\noalign{}
维度 & 详情 \\
\midrule\noalign{}
\endhead
\bottomrule\noalign{}
\endlastfoot
\textbf{成立时间} & 2018年 \\
\textbf{在线GPU数} & \textgreater20,000张（2026年） \\
\textbf{覆盖区域} & 北美、欧洲、亚太40+数据中心 \\
\textbf{供给方准入门槛} & 极低（个人即可注册为Host） \\
\textbf{平台抽佣} & ≈15\%（从Host收入中扣除） \\
\textbf{核心优势} & GPU种类最全、流动性最好、价格最低 \\
\textbf{核心劣势} & 可靠性无SLA、数据安全存疑、客服响应慢 \\
\textbf{2024年关键动向} & SOC
2认证、Serverless产品线、CLI/SDK全面升级 \\
\end{longtable}
}

\paragraph{6.1.2 RunPod ---
面向开发者的GPU云}\label{runpod-ux9762ux5411ux5f00ux53d1ux8005ux7684gpuux4e91}

{\def\LTcaptype{none} % do not increment counter
\begin{longtable}[]{@{}ll@{}}
\toprule\noalign{}
维度 & 详情 \\
\midrule\noalign{}
\endhead
\bottomrule\noalign{}
\endlastfoot
\textbf{成立时间} & 2022年 \\
\textbf{核心产品} & Pods（按需GPU）+ Serverless（自动扩缩容） \\
\textbf{GPU类型} & B200到RTX 4090，30+SKU \\
\textbf{全球区域} & 8+区域（含欧洲、亚太） \\
\textbf{关键差异化} & 热启动（Hot Start）Serverless，消除冷启动延迟 \\
\textbf{供给端策略} & 自有+合作数据中心为主，个人Host为辅 \\
\textbf{平台抽佣} & 约20-25\%（隐含在定价中） \\
\textbf{适用场景} & 生产级AI推理、Agent部署 \\
\end{longtable}
}

\paragraph{6.1.3 Akash Network ---
去中心化云计算市场}\label{akash-network-ux53bbux4e2dux5fc3ux5316ux4e91ux8ba1ux7b97ux5e02ux573a}

{\def\LTcaptype{none} % do not increment counter
\begin{longtable}[]{@{}ll@{}}
\toprule\noalign{}
维度 & 详情 \\
\midrule\noalign{}
\endhead
\bottomrule\noalign{}
\endlastfoot
\textbf{成立时间} & 2018年（主网2020年上线） \\
\textbf{底层技术} & Cosmos SDK区块链 + AKT代币 \\
\textbf{核心机制} & 反向拍卖（租户出价，提供商竞价） \\
\textbf{GPU定价} & H200 ≈\$2.59/hr, A100 ≈\$0.80/hr（2026年数据） \\
\textbf{关键生态} & AkashML（专注AI推理）、Razer合作案例 \\
\textbf{独特优势} & 抗审查、无需KYC、代币经济激励 \\
\textbf{主要局限} & 代币波动风险、技术门槛高、生态较小 \\
\end{longtable}
}

\paragraph{6.1.4
其他重要玩家}\label{ux5176ux4ed6ux91cdux8981ux73a9ux5bb6}

{\def\LTcaptype{none} % do not increment counter
\begin{longtable}[]{@{}
  >{\raggedright\arraybackslash}p{(\linewidth - 4\tabcolsep) * \real{0.2857}}
  >{\raggedright\arraybackslash}p{(\linewidth - 4\tabcolsep) * \real{0.2857}}
  >{\raggedright\arraybackslash}p{(\linewidth - 4\tabcolsep) * \real{0.4286}}@{}}
\toprule\noalign{}
\begin{minipage}[b]{\linewidth}\raggedright
平台
\end{minipage} & \begin{minipage}[b]{\linewidth}\raggedright
定位
\end{minipage} & \begin{minipage}[b]{\linewidth}\raggedright
关键数据
\end{minipage} \\
\midrule\noalign{}
\endhead
\bottomrule\noalign{}
\endlastfoot
\textbf{Render Network} & 去中心化3D渲染 &
RNDR代币，专注Octane/Redshift渲染 \\
\textbf{io.net} & Solana去中心化GPU网络 & 2024年推出，AI训练+推理 \\
\textbf{Salad} & 消费级PC共享 & 190+国家，以游戏奖励/礼品卡结算 \\
\textbf{Golem Network} & 通用去中心化算力 &
最早项目之一（2016年），以太坊生态 \\
\textbf{CoreWeave} & 从矿场转型GPU云 &
2024年估值\$190亿，已脱离P2P范畴 \\
\textbf{FluidStack} & 算力聚合+B2B & 汇集个人节点，面向企业客户 \\
\end{longtable}
}

\subsubsection{6.2
中国市场竞争格局}\label{ux4e2dux56fdux5e02ux573aux7adeux4e89ux683cux5c40}

{\def\LTcaptype{none} % do not increment counter
\begin{longtable}[]{@{}
  >{\raggedright\arraybackslash}p{(\linewidth - 6\tabcolsep) * \real{0.2222}}
  >{\raggedright\arraybackslash}p{(\linewidth - 6\tabcolsep) * \real{0.2222}}
  >{\raggedright\arraybackslash}p{(\linewidth - 6\tabcolsep) * \real{0.2222}}
  >{\raggedright\arraybackslash}p{(\linewidth - 6\tabcolsep) * \real{0.3333}}@{}}
\toprule\noalign{}
\begin{minipage}[b]{\linewidth}\raggedright
平台
\end{minipage} & \begin{minipage}[b]{\linewidth}\raggedright
类型
\end{minipage} & \begin{minipage}[b]{\linewidth}\raggedright
特点
\end{minipage} & \begin{minipage}[b]{\linewidth}\raggedright
合规状态
\end{minipage} \\
\midrule\noalign{}
\endhead
\bottomrule\noalign{}
\endlastfoot
\textbf{趋动云 (VirtAI)} & GPU池化平台 & 猎户座Orion软件，面向企业 & ✅
合规 \\
\textbf{矩池云 (Matpool)} & 共享GPU租赁 & 面向AI开发者，按时租赁 & ✅
合规 \\
\textbf{阿里云GPU共享} & 云厂商延伸 & cGPU/MIG分片共享 & ✅
合规（企业级） \\
\textbf{极客云/极链云} & P2P算力共享 & 个人可出租，规模较小 & ⚠️
模糊地带 \\
\textbf{AutoDL} & AI开发平台 & GPU租赁+Jupyter环境 & ✅ 合规 \\
\textbf{滴滴云（已退出）} & 公有云 & 2024年关停GPU租赁 & N/A \\
\end{longtable}
}

\textbf{中国市场关键判断：} -
真正的C2C个人算力共享在中国几乎处于空白状态，最大障碍是合规不确定性 -
现有平台多为B2C模式（平台自持或合作数据中心硬件），个人出租通道极少 -
矿场转型算力中心是更实际的路径（但面临电力、环评等审批） -
``东数西算''战略下，西部算力中心是大方向，个人节点的角色有待政策明确

\subsubsection{6.3
竞争要素总结}\label{ux7adeux4e89ux8981ux7d20ux603bux7ed3}

\begin{verbatim}
市场竞争核心维度评估：

价格竞争力  ████████████░░░░░░  Vast.ai > Akash > RunPod
供给丰富度  ██████████████░░░░  Vast.ai > Salad > RunPod
可靠性      ████████████████░░  RunPod > Vast.ai > Akash
易用性      ██████████████░░░░  RunPod > Salad > Vast.ai
数据安全    ████████████████░░  RunPod > Vast.ai > 区块链类
合规度      ██████████████████  传统云 > P2P平台 > 区块链类
\end{verbatim}

\begin{center}\rule{0.5\linewidth}{0.5pt}\end{center}

\pagebreak

\subsection{7.
个人出租收益模型深度分析}\label{ux4e2aux4ebaux51faux79dfux6536ux76caux6a21ux578bux6df1ux5ea6ux5206ux6790}

\subsubsection{7.1
主流GPU出租收益测算}\label{ux4e3bux6d41gpuux51faux79dfux6536ux76caux6d4bux7b97}

以下测算基于2026年6月Vast.ai和RunPod平台实际价格数据，假设平均出租率60\%（月出租约430小时）：

{\def\LTcaptype{none} % do not increment counter
\begin{longtable}[]{@{}
  >{\raggedright\arraybackslash}p{(\linewidth - 16\tabcolsep) * \real{0.0928}}
  >{\raggedright\arraybackslash}p{(\linewidth - 16\tabcolsep) * \real{0.1134}}
  >{\raggedright\arraybackslash}p{(\linewidth - 16\tabcolsep) * \real{0.1134}}
  >{\raggedright\arraybackslash}p{(\linewidth - 16\tabcolsep) * \real{0.1649}}
  >{\raggedright\arraybackslash}p{(\linewidth - 16\tabcolsep) * \real{0.1237}}
  >{\raggedright\arraybackslash}p{(\linewidth - 16\tabcolsep) * \real{0.1134}}
  >{\raggedright\arraybackslash}p{(\linewidth - 16\tabcolsep) * \real{0.0928}}
  >{\raggedright\arraybackslash}p{(\linewidth - 16\tabcolsep) * \real{0.0928}}
  >{\raggedright\arraybackslash}p{(\linewidth - 16\tabcolsep) * \real{0.0928}}@{}}
\toprule\noalign{}
\begin{minipage}[b]{\linewidth}\raggedright
GPU型号
\end{minipage} & \begin{minipage}[b]{\linewidth}\raggedright
平均时租价
\end{minipage} & \begin{minipage}[b]{\linewidth}\raggedright
月理论收入
\end{minipage} & \begin{minipage}[b]{\linewidth}\raggedright
月实际收入(60\%)
\end{minipage} & \begin{minipage}[b]{\linewidth}\raggedright
平台费(15\%)
\end{minipage} & \begin{minipage}[b]{\linewidth}\raggedright
月电力成本
\end{minipage} & \begin{minipage}[b]{\linewidth}\raggedright
月净收益
\end{minipage} & \begin{minipage}[b]{\linewidth}\raggedright
硬件价格
\end{minipage} & \begin{minipage}[b]{\linewidth}\raggedright
回本周期
\end{minipage} \\
\midrule\noalign{}
\endhead
\bottomrule\noalign{}
\endlastfoot
\textbf{RTX 4090 24GB} & \$0.35-0.55 & \$252-396 & \$151-238 & -\$23-36
& -\$20-35 & \textbf{\$108-167} & \$1,600 & 10-15月 \\
\textbf{RTX 3090 24GB} & \$0.20-0.35 & \$144-252 & \$86-151 & -\$13-23 &
-\$18-30 & \textbf{\$55-98} & \$800 & 8-15月 \\
\textbf{RTX 4080 16GB} & \$0.25-0.40 & \$180-288 & \$108-173 & -\$16-26
& -\$15-25 & \textbf{\$77-122} & \$1,000 & 8-13月 \\
\textbf{RTX 4070 Ti} & \$0.18-0.30 & \$130-216 & \$78-130 & -\$12-20 &
-\$12-18 & \textbf{\$54-92} & \$750 & 8-14月 \\
\textbf{A100 80GB} & \$0.80-1.50 & \$576-1,080 & \$346-648 & -\$52-97 &
-\$30-50 & \textbf{\$264-501} & \$10,000 & 20-38月 \\
\textbf{RTX 6000 Ada} & \$0.60-1.00 & \$432-720 & \$259-432 & -\$39-65 &
-\$20-30 & \textbf{\$200-337} & \$6,800 & 20-34月 \\
\textbf{GTX 1660 Ti} & \$0.05-0.10 & \$36-72 & \$22-43 & -\$3-6 &
-\$8-12 & \textbf{\$11-25} & \$200 & 8-18月 \\
\end{longtable}
}

\subsubsection{7.2
影响收益的关键变量}\label{ux5f71ux54cdux6536ux76caux7684ux5173ux952eux53d8ux91cf}

\paragraph{7.2.1
出租率（最关键变量）}\label{ux51faux79dfux7387ux6700ux5173ux952eux53d8ux91cf}

{\def\LTcaptype{none} % do not increment counter
\begin{longtable}[]{@{}lll@{}}
\toprule\noalign{}
出租率 & 驱动因素 & 典型场景 \\
\midrule\noalign{}
\endhead
\bottomrule\noalign{}
\endlastfoot
80-95\% & 低价策略（Bottom 20\%价格区间） & 低端GPU、高竞争区 \\
60-80\% & 定价略低于市场中位价 & 主流GPU、热点区域 \\
40-60\% & 定价适中、新Host信誉低 & 新注册Host \\
\textless30\% & 定价过高、配置冷门、频繁下线 & 定价策略失误 \\
\end{longtable}
}

\paragraph{7.2.2
电力成本（地域差异显著）}\label{ux7535ux529bux6210ux672cux5730ux57dfux5deeux5f02ux663eux8457}

{\def\LTcaptype{none} % do not increment counter
\begin{longtable}[]{@{}llll@{}}
\toprule\noalign{}
地区 & 居民电价 & 月电力成本(RTX 4090) & 对净收益影响 \\
\midrule\noalign{}
\endhead
\bottomrule\noalign{}
\endlastfoot
北美（平均） & \$0.12/kWh & ≈\$22 & 基准 \\
欧洲（德国） & \$0.35/kWh & ≈\$63 & 严重侵蚀利润 \\
中国（居民） & ¥0.5-0.6/kWh & ≈¥90-110 (\$12-15) & 成本优势明显 \\
中国（商业） & ¥1.0-1.2/kWh & ≈¥180-220 (\$25-30) & 与北美持平 \\
\end{longtable}
}

\paragraph{7.2.3
影响出租率的核心因素}\label{ux5f71ux54cdux51faux79dfux7387ux7684ux6838ux5fc3ux56e0ux7d20}

\begin{enumerate}
\def\labelenumi{\arabic{enumi}.}
\tightlist
\item
  \textbf{定价策略：} 低于市场中位价10-20\%可显著提升出租率
\item
  \textbf{信誉积累：} 累计出租时长、用户评价、在线稳定性
\item
  \textbf{地理位置：}
  北美\textgreater 欧洲\textgreater 亚太（需求分布不均）
\item
  \textbf{网络质量：} 上行带宽\textgreater100Mbps，延迟\textless50ms为佳
\item
  \textbf{硬件可靠性：} 频繁掉线会被平台降权
\item
  \textbf{GPU型号热度：} RTX 4090/A100/H100常年供不应求
\end{enumerate}

\subsubsection{7.3
多卡规模效应}\label{ux591aux5361ux89c4ux6a21ux6548ux5e94}

{\def\LTcaptype{none} % do not increment counter
\begin{longtable}[]{@{}
  >{\raggedright\arraybackslash}p{(\linewidth - 8\tabcolsep) * \real{0.1395}}
  >{\raggedright\arraybackslash}p{(\linewidth - 8\tabcolsep) * \real{0.1860}}
  >{\raggedright\arraybackslash}p{(\linewidth - 8\tabcolsep) * \real{0.2558}}
  >{\raggedright\arraybackslash}p{(\linewidth - 8\tabcolsep) * \real{0.2093}}
  >{\raggedright\arraybackslash}p{(\linewidth - 8\tabcolsep) * \real{0.2093}}@{}}
\toprule\noalign{}
\begin{minipage}[b]{\linewidth}\raggedright
规模
\end{minipage} & \begin{minipage}[b]{\linewidth}\raggedright
GPU数量
\end{minipage} & \begin{minipage}[b]{\linewidth}\raggedright
月均净收益
\end{minipage} & \begin{minipage}[b]{\linewidth}\raggedright
边际成本
\end{minipage} & \begin{minipage}[b]{\linewidth}\raggedright
关键挑战
\end{minipage} \\
\midrule\noalign{}
\endhead
\bottomrule\noalign{}
\endlastfoot
个人入门 & 1-2张 & \$100-350 & 低（利用现有PC） & 家庭网络/电力限制 \\
个人进阶 & 3-6张 & \$400-1,200 & 中（需专用机架） &
散热、噪音、电力扩容 \\
小型矿场转型 & 10-50张 & \$1,500-8,000 & 中（场地+运维） &
商业电价、网络带宽 \\
专业算力商 & 50-200张 & \$8,000-30,000 & 高（运维团队） &
资本投入、合规 \\
\end{longtable}
}

\begin{center}\rule{0.5\linewidth}{0.5pt}\end{center}

\pagebreak

\subsection{8.
政策与法规环境}\label{ux653fux7b56ux4e0eux6cd5ux89c4ux73afux5883}

\subsubsection{8.1
全球监管态势}\label{ux5168ux7403ux76d1ux7ba1ux6001ux52bf}

\paragraph{8.1.1 美国}\label{ux7f8eux56fd}

\begin{itemize}
\tightlist
\item
  \textbf{AI芯片出口管制：}
  2022年10月、2023年10月、2025年1月连续升级对华芯片出口限制，H100/B200/A100等高端GPU对华出口需BIS许可证
\item
  \textbf{算力即服务监管：} 2024年起要求IaaS提供商实施''Know Your
  Customer''(KYC)政策，防止受制裁实体获取AI算力
\item
  \textbf{税务合规：}
  Host收入需自行申报，平台逐步提供1099税表（年收入\textgreater\$600）
\end{itemize}

\paragraph{8.1.2 欧盟}\label{ux6b27ux76df}

\begin{itemize}
\tightlist
\item
  \textbf{AI Act (2024年)：}
  对通用AI(GPAI)模型提供者的算力使用有透明度要求
\item
  \textbf{GDPR延伸：}
  在第三方GPU上处理个人数据需满足数据保护要求，数据跨境传输受限
\item
  \textbf{能源效率指令：} 数据中心能效要求可能间接影响大规模算力出租
\end{itemize}

\paragraph{8.1.3
新加坡/东南亚}\label{ux65b0ux52a0ux5761ux4e1cux5357ux4e9a}

\begin{itemize}
\tightlist
\item
  相对开放的监管环境，多个平台选择新加坡作为亚太枢纽
\item
  加密货币友好的政策使其成为Web3算力平台的重要节点
\end{itemize}

\subsubsection{8.2
中国政策环境（重点分析）}\label{ux4e2dux56fdux653fux7b56ux73afux5883ux91cdux70b9ux5206ux6790}

\paragraph{8.2.1
算力基础设施政策}\label{ux7b97ux529bux57faux7840ux8bbeux65bdux653fux7b56}

{\def\LTcaptype{none} % do not increment counter
\begin{longtable}[]{@{}
  >{\raggedright\arraybackslash}p{(\linewidth - 6\tabcolsep) * \real{0.2292}}
  >{\raggedright\arraybackslash}p{(\linewidth - 6\tabcolsep) * \real{0.1875}}
  >{\raggedright\arraybackslash}p{(\linewidth - 6\tabcolsep) * \real{0.1875}}
  >{\raggedright\arraybackslash}p{(\linewidth - 6\tabcolsep) * \real{0.3958}}@{}}
\toprule\noalign{}
\begin{minipage}[b]{\linewidth}\raggedright
政策/文件
\end{minipage} & \begin{minipage}[b]{\linewidth}\raggedright
发布时间
\end{minipage} & \begin{minipage}[b]{\linewidth}\raggedright
核心内容
\end{minipage} & \begin{minipage}[b]{\linewidth}\raggedright
对个人算力出租的影响
\end{minipage} \\
\midrule\noalign{}
\endhead
\bottomrule\noalign{}
\endlastfoot
\textbf{``东数西算''工程} & 2022年2月 &
8大枢纽+10大数据中心集群，总投资\textgreater¥4,000亿 &
推动算力网络化，长期利好 \\
\textbf{《算力基础设施高质量发展行动计划》} & 2023年10月 &
2025年算力规模\textgreater300 EFLOPS & 算力纳入基础设施，规范化发展 \\
\textbf{《关于深入实施''东数西算''工程加快构建全国一体化算力网的实施意见》}
& 2023年12月 & 算力调度平台、算力交易机制 &
首次提及''算力交易''，为个人节点参与打开窗口 \\
\textbf{《数据出境安全评估办法》} & 2022年9月 & 数据出境需安全评估 &
限制跨境算力共享中的数据流动 \\
\end{longtable}
}

\paragraph{8.2.2
加密货币相关限制}\label{ux52a0ux5bc6ux8d27ux5e01ux76f8ux5173ux9650ux5236}

\begin{itemize}
\tightlist
\item
  \textbf{2021年9月：} 十部委联合发文，虚拟货币''挖矿''列为淘汰类产业
\item
  \textbf{影响：} 基于代币激励的区块链算力平台（Akash、Render
  Network）在中国面临本质合规障碍
\item
  \textbf{区别：}
  以法币结算的P2P算力出租不属于''挖矿''，但灰色地带仍然存在
\end{itemize}

\paragraph{8.2.3
合规路径分析}\label{ux5408ux89c4ux8defux5f84ux5206ux6790}

\begin{verbatim}
中国个人参与算力出租的可能路径：

✅ 相对合规路径：
  1. 通过国内合规平台（如矩池云、AutoDL等）出租——但目前对个人开放有限
  2. 出租给国内企业/机构——数据不出境，合规风险较低
  3. 作为"东数西算"节点——长期方向，需对接算力调度平台

⚠️ 灰色地带：
  1. 在Vast.ai等海外平台出租给境外用户——可能涉及数据跨境问题
  2. 家用宽带商业使用——违反运营商ToS
  3. 未办理ICP/增值电信许可证——个人出租算力是否属于电信业务存争议

❌ 高风险行为：
  1. 参与加密货币结算的算力平台——触犯虚拟货币禁令
  2. 明知租用方用于违法用途（如生成深度伪造）而不阻止
  3. 大规模商用但未进行税务登记和申报
\end{verbatim}

\subsubsection{8.3
知识产权与法律责任}\label{ux77e5ux8bc6ux4ea7ux6743ux4e0eux6cd5ux5f8bux8d23ux4efb}

\begin{itemize}
\tightlist
\item
  \textbf{模型权重泄露风险：}
  Host有权访问运行的容器，技术上可提取租用方的模型权重
\item
  \textbf{生成内容合法性：}
  GPU持有者是否对租用方生成的违法内容负责------目前法律空白
\item
  \textbf{AI生成内容版权：} 使用他人算力生成的AI内容版权归属尚无定论
\end{itemize}

\begin{center}\rule{0.5\linewidth}{0.5pt}\end{center}

\pagebreak

\subsection{9.
技术趋势与创新前沿}\label{ux6280ux672fux8d8bux52bfux4e0eux521bux65b0ux524dux6cbf}

\subsubsection{9.1
关键技术演进}\label{ux5173ux952eux6280ux672fux6f14ux8fdb}

\paragraph{9.1.1
GPU虚拟化与隔离}\label{gpuux865aux62dfux5316ux4e0eux9694ux79bb}

{\def\LTcaptype{none} % do not increment counter
\begin{longtable}[]{@{}
  >{\raggedright\arraybackslash}p{(\linewidth - 4\tabcolsep) * \real{0.3000}}
  >{\raggedright\arraybackslash}p{(\linewidth - 4\tabcolsep) * \real{0.4000}}
  >{\raggedright\arraybackslash}p{(\linewidth - 4\tabcolsep) * \real{0.3000}}@{}}
\toprule\noalign{}
\begin{minipage}[b]{\linewidth}\raggedright
技术
\end{minipage} & \begin{minipage}[b]{\linewidth}\raggedright
成熟度
\end{minipage} & \begin{minipage}[b]{\linewidth}\raggedright
说明
\end{minipage} \\
\midrule\noalign{}
\endhead
\bottomrule\noalign{}
\endlastfoot
\textbf{NVIDIA MIG} & 成熟（A100+/H100） &
硬件级GPU分区，单卡可切分为7个独立实例 \\
\textbf{NVIDIA vGPU} & 成熟 & 虚拟机GPU共享，授权费用较高 \\
\textbf{MIG-on-RTX} & 实验阶段 & RTX 6000
Ada已支持MIG，消费级卡尚不支持 \\
\textbf{GPU池化（rCUDA等）} & 早期 & 跨网络GPU调用，延迟是瓶颈 \\
\end{longtable}
}

\paragraph{9.1.2
可信执行环境（TEE）}\label{ux53efux4fe1ux6267ux884cux73afux5883tee}

\begin{itemize}
\tightlist
\item
  \textbf{NVIDIA Confidential
  Computing}（H100+支持）：硬件级加密，Host无法读取租用方数据------可能彻底改变P2P算力的安全范式
\item
  \textbf{AMD SEV-SNP}：AMD GPU路线图上的机密计算方案
\item
  \textbf{Intel TDX}：面向数据中心场景
\end{itemize}

\textbf{关键判断：}
TEE的普及将是个人算力出租从''低端外溢需求''走向''企业级可用''的分水岭。

\paragraph{9.1.3
分布式训练优化}\label{ux5206ux5e03ux5f0fux8badux7ec3ux4f18ux5316}

\begin{itemize}
\tightlist
\item
  \textbf{Petals、FlexGen等分布式推理/训练框架}：允许在异构、非稳定网络上进行大模型协作
\item
  \textbf{DiLoCo（Google
  DeepMind）}：分布式低通信优化，降低对网络互联的要求
\item
  \textbf{趋势：}
  算法优化正在降低分布式训练的通信带宽需求，使P2P网络上的协作训练更加可行
\end{itemize}

\subsubsection{9.2
新兴商业模式}\label{ux65b0ux5174ux5546ux4e1aux6a21ux5f0f}

\paragraph{9.2.1 算力金融化}\label{ux7b97ux529bux91d1ux878dux5316}

\begin{itemize}
\tightlist
\item
  \textbf{算力期货/期权：} 提前锁定未来算力价格和可用性
\item
  \textbf{GPU算力NFT：} 代币化算力权益（如Render Network的RNDR）
\item
  \textbf{算力Staking/借贷：} 类似DeFi的算力金融产品
\end{itemize}

\paragraph{9.2.2 AI-Fi（AI + DeFi融合）}\label{ai-fiai-defiux878dux5408}

\begin{itemize}
\tightlist
\item
  去中心化AI模型训练市场：贡献算力训练公共模型，按贡献分配代币奖励
\item
  模型推理收益共享：运行AI推理节点获得模型使用费分成
\end{itemize}

\paragraph{9.2.3 边缘-云协同}\label{ux8fb9ux7f18-ux4e91ux534fux540c}

\begin{itemize}
\tightlist
\item
  个人节点处理低延迟推理（游戏AI、AR/VR），云端处理大模型训练
\item
  5G MEC（移动边缘计算）+ 家庭GPU节点的分层调度
\end{itemize}

\begin{center}\rule{0.5\linewidth}{0.5pt}\end{center}

\pagebreak

\subsection{10.
战略机遇评估}\label{ux6218ux7565ux673aux9047ux8bc4ux4f30}

\subsubsection{10.1 机遇矩阵}\label{ux673aux9047ux77e9ux9635}

\begin{verbatim}
                    高
                    ↑
  竞争强度          │   ╔═══════════════╦═══════════════╗
                    │   ║  ② 算力聚合   ║  ① 合规算力   ║
                    │   ║  运营商      ║  交易平台     ║
                    │   ║  (中短期)    ║  (中长期)     ║
                    │   ╠═══════════════╬═══════════════╣
                    │   ║  ③ 纯粹出租  ║  ④ TEE安全   ║
                    │   ║  个人GPU     ║  算力方案     ║
                    │   ║  (低门槛)    ║  (前沿)       ║
                    │   ╚═══════════════╩═══════════════╝
                   低
                    低  ←──────────────────────→  高
                              市场空间/增长潜力
\end{verbatim}

\subsubsection{10.2
四大赛道深度评估}\label{ux56dbux5927ux8d5bux9053ux6df1ux5ea6ux8bc4ux4f30}

\paragraph{赛道①：合规化算力交易平台（★★★★★）}\label{ux8d5bux9053ux2460ux5408ux89c4ux5316ux7b97ux529bux4ea4ux6613ux5e73ux53f0}

\begin{itemize}
\tightlist
\item
  \textbf{市场空白：} 中国尚无合规的C2C算力交易平台
\item
  \textbf{政策窗口：} ``东数西算''提及''算力交易机制''，政策信号积极
\item
  \textbf{核心壁垒：}
  牌照（增值电信业务许可证）、数据合规体系、银行级结算
\item
  \textbf{目标用户：} 国内AI中小企业和个人开发者
\item
  \textbf{进入建议：} 与地方政府合作，在''东数西算''枢纽节点试点
\end{itemize}

\paragraph{赛道②：算力聚合运营商（★★★★☆）}\label{ux8d5bux9053ux2461ux7b97ux529bux805aux5408ux8fd0ux8425ux5546}

\begin{itemize}
\tightlist
\item
  \textbf{模式：} 汇集个人/小型GPU节点，统一标准化后向B端客户销售
\item
  \textbf{价值：} 解决一致性问题（SLA、安全、运维），抽成高于普通Host
\item
  \textbf{代表：} FluidStack（海外），国内空白
\item
  \textbf{核心能力：} 运维自动化、GPU测试认证、7×24监控
\item
  \textbf{盈利模型：}
  从Host端\$0.30/hr采购，到Client端\$0.60/hr出售，毛利率约40-50\%
\end{itemize}

\paragraph{赛道③：个人GPU出租（★★★☆☆）}\label{ux8d5bux9053ux2462ux4e2aux4ebagpuux51faux79df}

\begin{itemize}
\tightlist
\item
  \textbf{适合人群：} 已持有高端GPU的游戏玩家、内容创作者
\item
  \textbf{核心策略：} 平台组合（Vast.ai + RunPod）+ 时段优化 + 低价竞争
\item
  \textbf{关键风险：} 硬件损耗、电力成本上升、价格竞争加剧
\item
  \textbf{2026年判断：} 黄金期已过，边际收益递减，不建议新购卡专门出租
\end{itemize}

\paragraph{赛道④：TEE安全算力方案（★★☆☆☆
前沿）}\label{ux8d5bux9053ux2463teeux5b89ux5168ux7b97ux529bux65b9ux6848-ux524dux6cbf}

\begin{itemize}
\tightlist
\item
  \textbf{技术前提：} NVIDIA Confidential Computing下沉到消费级GPU
\item
  \textbf{时间窗口：} 预计3-5年后
\item
  \textbf{战略价值：} 一旦TEE普及，个人算力可进入金融、医疗等敏感场景
\end{itemize}

\begin{center}\rule{0.5\linewidth}{0.5pt}\end{center}

\pagebreak

\subsection{11.
进入策略与行动路线图}\label{ux8fdbux5165ux7b56ux7565ux4e0eux884cux52a8ux8defux7ebfux56fe}

\subsubsection{11.1
面向个人的行动路线}\label{ux9762ux5411ux4e2aux4ebaux7684ux884cux52a8ux8defux7ebf}

\begin{verbatim}
Phase 1: 试水验证（第1-2月）
├── 注册Vast.ai/RunPod Host账号
├── 完成KYC验证
├── 部署1张现有GPU（不需额外采购）
├── 定价策略：低于市场中位价15%（快速积累信誉）
├── 目标：30天出租率>50%，验证收益可行性

Phase 2: 优化扩展（第3-6月）
├── 根据数据调整定价和平台组合
├── 优化散热（机架+空调或专用通风）
├── 升级网络（有线连接+备用网络）
├── 如出租率>70%，可考虑加卡
├── 目标：月净收益>$200/卡，出租率>65%

Phase 3: 规模化（第7-12月）
├── 建立标准化运维流程（监控+报警+自动重启）
├── 探索多平台部署（对冲单一平台风险）
├── 考虑商业电价申请
├── 目标：4-8卡规模，月净收益>$1,000
\end{verbatim}

\subsubsection{11.2
面向创业者的平台构建路线}\label{ux9762ux5411ux521bux4e1aux8005ux7684ux5e73ux53f0ux6784ux5efaux8defux7ebf}

\begin{verbatim}
Phase 1: MVP（第1-6月）                    | Phase 2: 合规建设（第7-12月）
├── 技术栈选型（K8s+Docker+计量计费）        | ├── 申请增值电信业务许可证（ICP）
├── 核心功能：GPU注册→测试→上架→匹配→结算   | ├── 数据安全合规体系搭建
├── 10-20个种子Host节点                      | ├── 第三方支付牌照/合作
├── 目标：验证C2C匹配效率和经济模型           | ├── 用户协议与免责条款
                                            | ├── 目标：获得运营许可
Phase 3: 增长（第13-24月）                   | Phase 4: 生态（第25月+）
├── 对接"东数西算"算力调度平台                | ├── 开放API供第三方集成
├── 引入企业级SLA保障（TEE节点）              | ├── 算力金融产品（期货/期权）
├── Host激励计划（电费补贴、硬件融资）         | ├── 区域算力交易所
├── 目标：1,000+ GPU在线，月GMV>$50万         | ├── 目标：成为中国合规算力交易基础设施
\end{verbatim}

\begin{center}\rule{0.5\linewidth}{0.5pt}\end{center}

\pagebreak

\subsection{12.
商业模式与财务模型}\label{ux5546ux4e1aux6a21ux5f0fux4e0eux8d22ux52a1ux6a21ux578b}

\subsubsection{12.1
平台方商业模式}\label{ux5e73ux53f0ux65b9ux5546ux4e1aux6a21ux5f0f}

\paragraph{模式A：抽佣制（Vast.ai模式）}\label{ux6a21ux5f0faux62bdux4f63ux5236vast.aiux6a21ux5f0f}

\begin{itemize}
\tightlist
\item
  \textbf{营收：} 交易额×15\%
\item
  \textbf{假设：} 10,000活跃GPU × 月均\$150/卡 = \$150万月GMV × 15\% =
  \$22.5万/月
\item
  \textbf{成本：} 匹配引擎、支付网关、客服
\item
  \textbf{优势：} 轻资产、可规模化
\item
  \textbf{劣势：} 单价低、需要海量节点
\end{itemize}

\paragraph{模式B：加价制（FluidStack模式）}\label{ux6a21ux5f0fbux52a0ux4ef7ux5236fluidstackux6a21ux5f0f}

\begin{itemize}
\tightlist
\item
  \textbf{营收：} 买入-卖出价差（约50-100\%）
\item
  \textbf{假设：} 1,000 GPU × 月均\$300采购 × 1.7倍出售 =
  \$51万月营收，毛利率≈\$21万
\item
  \textbf{成本：} 采购成本+运维+SLA保障
\item
  \textbf{优势：} 单GPU收入高、品牌溢价
\item
  \textbf{劣势：} 重运营、扩展慢
\end{itemize}

\paragraph{模式C：混合制（RunPod模式）}\label{ux6a21ux5f0fcux6df7ux5408ux5236runpodux6a21ux5f0f}

\begin{itemize}
\tightlist
\item
  自有节点+Host网络的混合供给
\item
  自有节点保证SLA，Host网络补充弹性
\item
  通过Serverless产品实现更高附加值
\end{itemize}

\subsubsection{12.2 个人出租者财务模型（5张RTX
4090规模）}\label{ux4e2aux4ebaux51faux79dfux8005ux8d22ux52a1ux6a21ux578b5ux5f20rtx-4090ux89c4ux6a21}

{\def\LTcaptype{none} % do not increment counter
\begin{longtable}[]{@{}
  >{\raggedright\arraybackslash}p{(\linewidth - 4\tabcolsep) * \real{0.3333}}
  >{\raggedright\arraybackslash}p{(\linewidth - 4\tabcolsep) * \real{0.3333}}
  >{\raggedright\arraybackslash}p{(\linewidth - 4\tabcolsep) * \real{0.3333}}@{}}
\toprule\noalign{}
\begin{minipage}[b]{\linewidth}\raggedright
项目
\end{minipage} & \begin{minipage}[b]{\linewidth}\raggedright
月度
\end{minipage} & \begin{minipage}[b]{\linewidth}\raggedright
年度
\end{minipage} \\
\midrule\noalign{}
\endhead
\bottomrule\noalign{}
\endlastfoot
\textbf{收入} & & \\
GPU出租收入（60\%利用率×5卡，按第7章单卡基准） & \$755-1,190 &
\$9,060-14,280 \\
\textbf{核心成本} & & \\
电力（\$0.12/kWh，5卡≈750-1,000kWh/月） & -\$90-150 & -\$1,080-1,800 \\
平台抽佣（15\%） & -\$113-179 & -\$1,356-2,148 \\
硬件折旧（5张×\$1,600÷36月） & -\$222 & -\$2,664 \\
\textbf{核心净收益} & \textbf{\$330-639} & \textbf{\$3,960-7,668} \\
\textbf{规模化管理附加成本} & & \\
网络带宽升级（商用宽带/多线接入） & -\$30-50 & -\$360-600 \\
运维监控工具+远程管理 & -\$20-30 & -\$240-360 \\
散热/机架折旧（初期投入\$500÷24月） & -\$21 & -\$252 \\
\textbf{综合净收益} & \textbf{\$259-538} & \textbf{\$3,108-6,456} \\
\end{longtable}
}

\begin{quote}
\textbf{注：} 核心净收益按第7章单卡基准直接放大（\$108-167×5 =
\$540-835→实测60\%利用率后\$330-639）。规模化管理附加成本为多卡部署的增量费用，单卡场景不涉及。与单卡相比，5卡规模可通过批量采购硬件（折扣5-10\%）、统一运维降低单位管理成本，实际综合净收益可能高于表中保守估算值。
\end{quote}

\begin{itemize}
\tightlist
\item
  \textbf{初始投资：} 5×\$1,600 + \$500(机架/散热) = \$8,500
\item
  \textbf{核心回本周期：} 约13-26个月（仅计核心成本）
\item
  \textbf{综合回本周期：} 约16-33个月（含全部附加成本）
\item
  \textbf{ROI（年化，综合净收益÷总投资）：} 约37-76\%
\end{itemize}

\begin{center}\rule{0.5\linewidth}{0.5pt}\end{center}

\pagebreak

\subsection{13.
风险识别与缓释方案}\label{ux98ceux9669ux8bc6ux522bux4e0eux7f13ux91caux65b9ux6848}

\subsubsection{13.1 风险全景}\label{ux98ceux9669ux5168ux666f}

{\def\LTcaptype{none} % do not increment counter
\begin{longtable}[]{@{}lllll@{}}
\toprule\noalign{}
风险类别 & 风险项 & 概率 & 影响 & 风险等级 \\
\midrule\noalign{}
\endhead
\bottomrule\noalign{}
\endlastfoot
\textbf{市场风险} & GPU价格战导致时租价持续下降 & 高 & 中 & 🔴 高 \\
\textbf{市场风险} & AI算力需求转向ASIC/专用芯片 & 中 & 高 & 🟡 中高 \\
\textbf{政策风险} & 中国加强算力跨境流动管制 & 高 & 高 & 🔴 高 \\
\textbf{政策风险} & 平台被要求强制KYC/持牌经营 & 中高 & 中 & 🟡 中高 \\
\textbf{技术风险} & Host节点数据泄露/模型被窃 & 中 & 高 & 🟡 中高 \\
\textbf{技术风险} & 硬件故障导致租用方损失索赔 & 低 & 高 & 🟢 中低 \\
\textbf{运营风险} & 电力供应不稳/电费涨价 & 中 & 中 & 🟡 中 \\
\textbf{竞争风险} & NVIDIA推出消费级GPU租赁限制 & 低 & 极高 & 🟡 中 \\
\textbf{合规风险} & 税务/工商对个人算力出租定性 & 中高 & 中 & 🟡 中高 \\
\end{longtable}
}

\subsubsection{13.2
分级缓释方案}\label{ux5206ux7ea7ux7f13ux91caux65b9ux6848}

\paragraph{红色风险（需立即应对）}\label{ux7ea2ux8272ux98ceux9669ux9700ux7acbux5373ux5e94ux5bf9}

{\def\LTcaptype{none} % do not increment counter
\begin{longtable}[]{@{}
  >{\raggedright\arraybackslash}p{(\linewidth - 2\tabcolsep) * \real{0.4000}}
  >{\raggedright\arraybackslash}p{(\linewidth - 2\tabcolsep) * \real{0.6000}}@{}}
\toprule\noalign{}
\begin{minipage}[b]{\linewidth}\raggedright
风险
\end{minipage} & \begin{minipage}[b]{\linewidth}\raggedright
缓释措施
\end{minipage} \\
\midrule\noalign{}
\endhead
\bottomrule\noalign{}
\endlastfoot
GPU价格战 &
①差异化配置（预装模型/环境）②转向企业长期租约③降低持有成本（利用夜间谷电） \\
中国政策收紧 &
①优先出租给国内用户②准备国内合规平台备选③避免数据中心级算力跨境 \\
\end{longtable}
}

\paragraph{黄色风险（需监控和预案）}\label{ux9ec4ux8272ux98ceux9669ux9700ux76d1ux63a7ux548cux9884ux6848}

{\def\LTcaptype{none} % do not increment counter
\begin{longtable}[]{@{}
  >{\raggedright\arraybackslash}p{(\linewidth - 2\tabcolsep) * \real{0.4000}}
  >{\raggedright\arraybackslash}p{(\linewidth - 2\tabcolsep) * \real{0.6000}}@{}}
\toprule\noalign{}
\begin{minipage}[b]{\linewidth}\raggedright
风险
\end{minipage} & \begin{minipage}[b]{\linewidth}\raggedright
缓释措施
\end{minipage} \\
\midrule\noalign{}
\endhead
\bottomrule\noalign{}
\endlastfoot
ASIC替代 & ①关注NVIDIA路线图②GPU推理仍有生态优势③多元化卡型降低集中度 \\
数据安全 & ①物理隔离敏感网络②定期固件更新③购买网络安全保险 \\
平台合规 & ①分散多平台②关注平台合规动态③保留交易记录备查 \\
电费上涨 & ①选址电价优势区域②能效优化（降频降压）③光伏+储能自供电 \\
\end{longtable}
}

\begin{center}\rule{0.5\linewidth}{0.5pt}\end{center}

\pagebreak

\subsection{14.
附录：关键数据表与参考文献}\label{ux9644ux5f55ux5173ux952eux6570ux636eux8868ux4e0eux53c2ux8003ux6587ux732e}

\subsubsection{14.1
主要平台速查表}\label{ux4e3bux8981ux5e73ux53f0ux901fux67e5ux8868}

{\def\LTcaptype{none} % do not increment counter
\begin{longtable}[]{@{}
  >{\raggedright\arraybackslash}p{(\linewidth - 10\tabcolsep) * \real{0.1224}}
  >{\raggedright\arraybackslash}p{(\linewidth - 10\tabcolsep) * \real{0.1224}}
  >{\raggedright\arraybackslash}p{(\linewidth - 10\tabcolsep) * \real{0.2245}}
  >{\raggedright\arraybackslash}p{(\linewidth - 10\tabcolsep) * \real{0.1837}}
  >{\raggedright\arraybackslash}p{(\linewidth - 10\tabcolsep) * \real{0.1224}}
  >{\raggedright\arraybackslash}p{(\linewidth - 10\tabcolsep) * \real{0.2245}}@{}}
\toprule\noalign{}
\begin{minipage}[b]{\linewidth}\raggedright
平台
\end{minipage} & \begin{minipage}[b]{\linewidth}\raggedright
网址
\end{minipage} & \begin{minipage}[b]{\linewidth}\raggedright
供给端类型
\end{minipage} & \begin{minipage}[b]{\linewidth}\raggedright
结算方式
\end{minipage} & \begin{minipage}[b]{\linewidth}\raggedright
抽佣
\end{minipage} & \begin{minipage}[b]{\linewidth}\raggedright
中国可用性
\end{minipage} \\
\midrule\noalign{}
\endhead
\bottomrule\noalign{}
\endlastfoot
Vast.ai & vast.ai & 个人/小型DC & 法币(信用卡) & ≈15\% & ⚠️
需翻墙+境外支付 \\
RunPod & runpod.io & 个人/DC & 法币 & ≈20-25\% & ⚠️ 需翻墙 \\
Akash Network & akash.network & 无许可 & AKT代币 & 市场化 & ❌
代币违禁 \\
Render Network & rendernetwork.com & 无许可 & RNDR代币 & 市场化 & ❌
代币违禁 \\
Salad & salad.com & 消费级PC & 积分/礼品卡 & 未公开 & ❌ 不可用 \\
矩池云 & matpool.com & 平台自持 & 法币(支付宝) & 隐含定价 & ✅ 合规 \\
AutoDL & autodl.com & 平台自持 & 法币 & 隐含定价 & ✅ 合规 \\
\end{longtable}
}

\subsubsection{14.2
GPU经济性对比}\label{gpuux7ecfux6d4eux6027ux5bf9ux6bd4}

{\def\LTcaptype{none} % do not increment counter
\begin{longtable}[]{@{}llllllll@{}}
\toprule\noalign{}
GPU & VRAM & 购入价 & 功率 & 月均收益 & 月均成本 & 月净利 & 回本月 \\
\midrule\noalign{}
\endhead
\bottomrule\noalign{}
\endlastfoot
RTX 4090 & 24GB & \$1,600 & 450W & \$200 & \$110 & \$90 & 18 \\
RTX 3090 & 24GB & \$800 & 350W & \$110 & \$70 & \$40 & 20 \\
RTX 4080S & 16GB & \$1,000 & 320W & \$140 & \$70 & \$70 & 14 \\
RTX 4070S & 12GB & \$600 & 220W & \$80 & \$45 & \$35 & 17 \\
A100 80GB & 80GB & \$10,000 & 400W & \$450 & \$160 & \$290 & 34 \\
H100 80GB & 80GB & \$25,000 & 700W & \$1,200 & \$300 & \$900 & 28 \\
\end{longtable}
}

\subsubsection{14.3 术语表}\label{ux672fux8bedux8868}

{\def\LTcaptype{none} % do not increment counter
\begin{longtable}[]{@{}
  >{\raggedright\arraybackslash}p{(\linewidth - 4\tabcolsep) * \real{0.3333}}
  >{\raggedright\arraybackslash}p{(\linewidth - 4\tabcolsep) * \real{0.3333}}
  >{\raggedright\arraybackslash}p{(\linewidth - 4\tabcolsep) * \real{0.3333}}@{}}
\toprule\noalign{}
\begin{minipage}[b]{\linewidth}\raggedright
术语
\end{minipage} & \begin{minipage}[b]{\linewidth}\raggedright
英文
\end{minipage} & \begin{minipage}[b]{\linewidth}\raggedright
定义
\end{minipage} \\
\midrule\noalign{}
\endhead
\bottomrule\noalign{}
\endlastfoot
算力共享 & Compute Sharing & 将闲置计算资源出租给他人的经济模式 \\
GPU池化 & GPU Pooling & 将多张物理GPU虚拟化为统一资源池 \\
TEE & Trusted Execution Environment &
硬件级可信执行环境，保护数据和代码不被主机访问 \\
MIG & Multi-Instance GPU & NVIDIA的GPU硬件分区技术 \\
算力调度 & Compute Orchestration & 将计算任务动态分配到最优算力节点 \\
反向拍卖 & Reverse Auction & 租户出价、多个提供商竞价的价格发现机制 \\
时区套利 & Timezone Arbitrage & 利用时区差异实现算力的24小时利用 \\
AI-Fi & AI + DeFi & 去中心化金融与人工智能的交叉领域 \\
\end{longtable}
}

\subsubsection{14.4
参考文献与数据来源}\label{ux53c2ux8003ux6587ux732eux4e0eux6570ux636eux6765ux6e90}

\begin{enumerate}
\def\labelenumi{\arabic{enumi}.}
\tightlist
\item
  \textbf{Grand View Research --- GPU Cloud Market Report 2024} ---
  全球GPU云市场规模和增长预测
\item
  \textbf{Fortune Business Insights --- GPU as a Service Market
  Analysis} --- GPUaaS市场结构分析
\item
  \textbf{Messari --- DePIN Sector Report 2024} ---
  去中心化物理基础设施网络行业报告
\item
  \textbf{Vast.ai Platform Data} ---
  平台公开的GPU供给和定价数据（通过API获取）
\item
  \textbf{RunPod State of AI Infrastructure Report} ---
  AI基础设施趋势和部署数据
\item
  \textbf{Akash Network Economics Paper} ---
  去中心化云计算经济模型白皮书
\item
  \textbf{中国信息通信研究院 --- 《中国算力发展指数白皮书（2024）》} ---
  中国算力基础设施发展数据
\item
  \textbf{国务院 --- 《``东数西算''工程实施方案》} ---
  国家算力网络战略规划
\item
  \textbf{NVIDIA --- Confidential Computing Technical Brief} ---
  GPU可信执行环境技术文档
\item
  \textbf{Tom's Hardware / Reddit r/gpumining / r/VastAI} ---
  社区GPU出租收益实测数据
\item
  \textbf{Steam Hardware Survey (2026年5月)} --- 全球GPU持有量分布数据
\item
  \textbf{BIS --- Semiconductor Export Control Rules (2023-2025)} ---
  美国对华芯片出口管制政策
\end{enumerate}

\begin{center}\rule{0.5\linewidth}{0.5pt}\end{center}

\emph{本报告由中普咨询调研部完成，所有数据和判断基于截至2026年6月的公开信息和行业调研。报告中的预测性信息存在不确定性，不构成投资建议。}

\textbf{报告编写：} 调研部/苏砚清 \textbf{数据审核：} 分析部/顾维钧
\textbf{情报验证：} 情报部/沈千帆 \textbf{合规审查：} 合规部/钟正则
\textbf{签批：} 总经办/CEO

\emph{密级：内部使用 \textbar{} 未经授权不得外传}
\end{document}
